% 第 44 章。数字与教学模型：checks/check_ch44.py；图：figsrc/fig_ch44.py。
% 史料：naturalrate_facts（取数脚本 data/fetch/naturalrate44.py）；美国与日本的失业率与消费物价：naturalrate_us_1959_1985、
% naturalrate_japan_1953_1986（同一脚本）；英国：boe_millennium_uk_1954_1986（data/fetch/boe_millennium_gb44.py）。
\chapter{自然率与附加预期的菲利普斯曲线}\label{ch:naturalrate}

\begin{chapterguide}
\guidequestion{菲利普斯“菜单”为什么会漂移？}
\begin{guidegoals}
\item 推导附加预期的菲利普斯曲线与加速论假说；
\item 解释菲尔普斯（1967）与弗里德曼（1968）为什么预言长期菲利普斯曲线是垂直的；
\item 用适应性预期模拟反通胀的产出代价；
\item 用 1970 年代的数据检验这一预言。
\end{guidegoals}
\guideprereq{\chref{ch:phillips}\enspace\nameref{ch:phillips}}
\end{chapterguide}

1967 年 12 月 29 日是星期五。那天晚上，弗里德曼（Milton Friedman，1912—2006）在华盛顿美国经济学会的第八十届年会上发表会长演说\cite{friedman1968,loc44}。演说登在次年 3 月的《美国经济评论》上，题为《货币政策的作用》\cite{friedman1968}。他开头说，高就业、物价稳定与快速增长是经济政策的主要目标，这一点大家意见一致\cite{friedman1968}。他当晚只讲其中的一件工具，货币政策\cite{friedman1968}。
% 据: f68_date, f68_issue, f68_title, f68_goals, f68_topic, loc_friedman

他先讲货币政策做不到的事，挑了两件。它钉不住利率，也钉不住失业率，都只能维持很短的时间\cite{friedman1968}。他说，第二件与当时流行的想法更加相左\cite{friedman1968}。人们普遍相信，货币增长会刺激就业。他于是问，货币当局为什么不能定一个目标，比如失业 3\%，失业低于目标就收紧，高于目标就放松\cite{friedman1968}？他的回答是，道理与利率一样，在于政策的即时后果与延迟后果不同\cite{friedman1968}。
% 据: f68_limits, f68_grain, f68_target, f68_reason

那几年美国的失业离 3\% 不远。按《总统经济报告》的统计，1967 年的失业率是 3.8\%，消费物价比上年涨 3.1\%\cite{erp2000}。1969 年失业降到 3.5\%，物价涨了 5.5\%\cite{erp2000}。萨缪尔森与索洛 1960 年画的菜单上，失业 3\% 对应物价每年涨约 4.5\%（\chref{ch:phillips}）。到了 1975 年，失业升到 8.5\%，物价仍涨 9.1\%\cite{erp2000}。1980 年，失业 7.1\%，物价涨 13.5\%\cite{erp2000}。这些点都不在那张菜单上（本书的解读）。菲利普斯“菜单”为什么会漂移？
% 据: b40_1967, b62_1967, b40_1969, b62_1969, b40_1975, b62_1975, b40_1980, b62_1980

\section{当时的世界}\label{sec:naturalrate-world}

九年以后，弗里德曼在诺贝尔奖的演讲里回顾了这段历史。演讲是 1976 年 12 月 13 日作的\cite{friedman1976}。他说，二战以后，经济学界对通胀与失业的关系的分析经过了两个阶段，正在进入第三个\cite{friedman1976}。第一个阶段接受了向下倾斜的菲利普斯曲线，它被普遍解释为一种因果关系，似乎给决策者提供了一种稳定的取舍\cite{friedman1976}。决策者可以选一个低的失业目标，接受相应的通胀，也可以选低通胀甚至通缩，忍受更高的失业\cite{friedman1976}。
% 据: f76_date, f76_stages, f76_stable, f76_choose

经济学家忙着从各国各时期的证据里估计这个关系，也研究通胀与失业各自的社会得失，好选出合适的取舍\cite{friedman1976}。可是新的证据与它不合\cite{friedman1976}。在各国政府都在追求“充分就业”的战后，与某一失业水平相伴的通胀率，在一国之内往往逐渐升高，各国之间也相差很大\cite{friedman1976}。高通胀与高失业同时出现，得了“滞胀”这个难听的名字\cite{friedman1976}。
% 据: f76_busied, f76_failed, f76_rise, f76_stagflation

弗里德曼 1967 年的演说也描述了当时的风气。他拿 1920 年代末与当时对照\cite{friedman1968}。那时，货币政策的主要任务是稳定物价、维护金本位\cite{friedman1968}。如今，促进充分就业被放在首位，防止通胀是一个持续的、却明显次要的目标\cite{friedman1968}。他担心钟摆又摆得太远，人们要货币政策去做它做不到的事，反倒妨碍它作出能作的贡献\cite{friedman1968}。
% 据: f68_twenties, f68_then, f68_today, f68_pendulum

谁让通胀升高？按弗里德曼的解释，各国政府并没有把高通胀当作公开宣布的政策来推行\cite{friedman1976}。高通胀是别的政策的后果，尤其是充分就业的政策与扩大政府开支的福利国家政策\cite{friedman1976}。他说，各国政府都宣称要稳定物价，一阵通胀之后，压住它的压力也很大\cite{friedman1976}。1962 年，肯尼迪把 4\% 的失业率定为美国的临时目标（\chref{ch:phillips}）。
% 据: f76_policies, f76_proclaim

尼克松执政时，美国政府直接管过物价。Gordon 与 King 估计，这些管制把物价水平压低了 1.5 个百分点，取消管制又把它推高 2.5 个百分点\cite{gordon1982}。
% 据: gk_controls

1974 年的紧缩是谁定的？莫迪利亚尼（Franco Modigliani，1918—2003）与 Papademos 在 1975 年的《布鲁金斯经济活动论文》上作了分析\cite{modigliani1975,loc44,gnd44}。他们写道，经济的现状大体要追溯到前几个季度的货币政策，特别是 1974 年年中的严厉紧缩\cite{modigliani1975}。按他们的说法，联邦储备 1974 年的政策混合了两个目标\cite{modigliani1975}。一是把 M1 的增长大致控制在经济顾问委员会 1974 年报告预示的 6\% 以内，二是不让利率降得太快\cite{modigliani1975}。政府为 1974 年定的目标是实际收入增长 1\%，物价上涨约 7\%，两者合起来要求货币收入增长约 8\%\cite{modigliani1975}。6\% 的货币增长远远不够\cite{modigliani1975}。
% 据: mp_squeeze, mp_blend, mp_m1target, mp_admin, mp_income8, mp_toosmall, loc_modigliani, gnd_modigliani

第二季度的紧缩让联邦基金利率从 3 月初的不到 9\% 升到 7 月初的 13.5\%\cite{modigliani1975}。他们认为，联邦储备容忍利率这样急升，反映了它不愿让货币总量偏离最初的目标\cite{modigliani1975}。决策依据的货币统计也不可靠，后来修订，2 月到 5 月 M1 的年增长率从 9.7\% 改成了 7.6\%\cite{modigliani1975}。他们不相信决策者有意要 9\% 的失业，这离政府宣布或国会认可的任何目标都太远\cite{modigliani1975}。他们认为，毛病出在以货币总量与“有序市场”为目标\cite{modigliani1975}。
% 据: mp_funds, mp_aversion, mp_revised, mp_nine, mp_orderly

英国的首相也改了口。弗里德曼在演讲里引了卡拉汉（James Callaghan，1912—2005）1976 年 9 月 28 日在工党大会上的话\cite{friedman1976,loc44}。卡拉汉说，过去以为靠减税与扩大政府开支就能走出衰退、增加就业，这个选项已经不存在了\cite{friedman1976}。他说，就算它从前起过作用，也只是靠给经济注入更大剂量的通胀，接着来的是更高的失业\cite{friedman1976}。弗里德曼说，这番话与战后几乎每个西方政府的政策正好相反\cite{friedman1976}。
% 据: f76_callaghan_date, f76_callaghan, f76_callaghan2, f76_callaghan3, f76_remarkable, loc_callaghan

\section{自然失业率}\label{sec:naturalrate-friedman}

弗里德曼的论证借用了维克塞尔（Knut Wicksell，1851—1926）\cite{friedman1968,loc44}。维克塞尔区分了自然利率与市场利率（\chref{ch:wicksell}）。弗里德曼在演说里说，货币当局只有靠通胀，才能让市场利率低于自然利率\cite{friedman1968}。人们普遍预期通胀以后，名义的自然利率本身也会上升，要压住市场利率，就得通胀得更快\cite{friedman1968}。他说，就业市场上有一个与此对应的东西\cite{friedman1968}。
% 据: f68_wicksell, f68_inflation, f68_fisher, f68_counterpart, loc_wicksell

在任何时候，都有一个与实际工资结构的均衡相容的失业水平\cite{friedman1968}。失业低于它，劳动供不应求，实际工资有上升的压力。失业高于它，实际工资有下降的压力\cite{friedman1968}。弗里德曼把这个水平叫作“自然失业率”。照他的定义，它是瓦尔拉斯一般均衡方程组算出来的失业水平\cite{friedman1968}。方程组里要放进劳动市场与商品市场的实际结构，包括市场的不完全、供求的随机变动、打听空缺与人手的成本、流动的成本等等\cite{friedman1968}。他在脚注里说，它不一定是失业人数等于空缺岗位数的那一点\cite{friedman1968}。
% 据: f68_equilibrium, f68_excess, f68_natural, f68_fn3

他特意说明，“自然”并不是说它一成不变\cite{friedman1968}。决定它的许多市场特征是人为的，是政策造成的\cite{friedman1968}。他举美国为例，法定最低工资、Walsh-Healey 法与 Davis-Bacon 法以及工会的力量，都让自然失业率比没有它们时高\cite{friedman1968}。职业介绍所的改进与空缺、人手信息的改进，会让它降低\cite{friedman1968}。他用“自然”二字，理由与维克塞尔相同，是要把实际的力量与货币的力量分开\cite{friedman1968}。
% 据: f68_immutable, f68_manmade, f68_minimum, f68_lower, f68_reason_natural

菲利普斯的曲线错在哪里？弗里德曼说，它有一个根本的缺陷，没有区分名义工资与实际工资，正像维克塞尔没有区分名义利率与实际利率\cite{friedman1968}。照他的读法，菲利普斯的文章写给一个人人都预期物价稳定、不管实际怎样都不改这个预期的世界\cite{friedman1968}。假如人人都预期物价每年涨 75\% 以上，工资就得涨这么多，实际工资才不变\cite{friedman1968}。他说，几年前的巴西人就是这样预期的\cite{friedman1968}。巴西把物价的涨幅压到每年约 45\% 时，失业一开始急剧上升，因为工资还按旧的预期在涨\cite{friedman1968}。
% 据: f68_defect, f68_world, f68_brazil75, f68_brazil_did, f68_brazil45

弗里德曼 1976 年回忆，他们几个人从一开始就怀疑稳定的菲利普斯曲线，理由主要在理论上，不在经验上\cite{friedman1976}。他们认为，对就业要紧的是工资能买到多少商品与服务，以美元、英镑或克朗计的工资本身并不要紧\cite{friedman1976}。低失业确实意味着实际工资有上涨的压力，可只要物价降得更多，名义工资降了，实际工资照样可以更高\cite{friedman1976}。
% 据: f76_sceptical, f76_realwages, f76_pressure

他在脚注里补充，物价的平均涨幅在一段时期里比较稳定，按名义工资画的菲利普斯曲线才稳定\cite{friedman1968}。不同时期、不同国家的曲线高低不同，平均涨幅越高，曲线的位置越高\cite{friedman1968}。这就是菜单漂移的原因：菜单画在某一种预期上，预期一变，菜单就挪了地方（本书的解读）。萨缪尔森与索洛 1960 年猜过，过去的经验会养成通胀的倾向，他们没有判断菜单在长期里往哪个方向移（\chref{ch:phillips}）。
% 据: f68_fn5, f68_level

\subsubsection{为什么守不住 3\%}

回到 3\% 的目标。设自然率高于 3\%，开头物价稳定，失业高于 3\%\cite{friedman1968}。货币当局加快货币增长，开支增加\cite{friedman1968}。人们原先预期物价稳定，价格与工资都按这个预期定好了，增加的收入起初多半变成产出与就业\cite{friedman1968}。产品的价格比工资涨得快，雇主付的实际工资降了。雇员却按旧的物价估算，以为实际工资涨了\cite{friedman1968}。按弗里德曼的说法，雇主眼里的实际工资下降，雇员眼里的实际工资上涨，正是这一点让就业增加\cite{friedman1968}。
% 据: f68_suppose, f68_expansion, f68_initial, f68_exante

雇员不久就会把物价的上涨算进去，要求更高的名义工资\cite{friedman1968}。实际工资回升，失业回到原来的水平\cite{friedman1968}。要把失业守在 3\%，货币当局只能让货币增长得更快\cite{friedman1968}。他的结论是，市场的失业率只有靠通胀，而且是加速的通胀，才能压在自然率以下\cite{friedman1968}。反过来，目标定在自然率以上，就会带来加速的通缩\cite{friedman1968}。
% 据: f68_anticipate, f68_return, f68_more, f68_accelerating, f68_deflation

那么，把目标定在自然率上呢？弗里德曼说，货币当局不知道自然率是多少，还没有办法准确地估计它，它本身也会不时变动\cite{friedman1968}。他的概括是，通胀与失业之间总有暂时的取舍，没有永久的取舍\cite{friedman1968}。暂时的取舍来自没有预料到的通胀，一般也就是上升的通胀\cite{friedman1968}。上升的通胀可以降低失业，高通胀不能\cite{friedman1968}。“暂时”有多久？他凭对历史证据的考察，个人判断最初的效果持续两到五年，就业完全调整到新的通胀率，大约要二十年\cite{friedman1968}。
% 据: f68_cannot_know, f68_change, f68_tradeoff, f68_unanticipated, f68_rising, f68_howlong, f68_judgment, f68_years, f68_decades

他把结论换了一种说法。货币当局控制的是名义量，直接控制的是它自己的负债\cite{friedman1968}。它可以用这种控制钉住一个名义量，如汇率、物价水平、名义国民收入或某种定义的货币数量，也可以钉住名义量的变化率\cite{friedman1968}。它不能用它钉住实际量，如实际利率、失业率、实际国民收入\cite{friedman1968}。他还说，上面的时长只适用于美国经历过的那种幅度，南美国家那样大得多的变化，调整要快得多\cite{friedman1968}。
% 据: f68_nominal, f68_peg_nominal, f68_real, f68_south

\chref{ch:phillips}讲过休谟的看法，增加的货币在物价追上来以前刺激了产业。弗里德曼补上了长期的那一半，物价与预期追上来以后，失业回到自然率，货币增长加快只留下更高的通胀（本书的概括）。\chref{ch:generaltheory}的争论问的是，工资与物价若能完全灵活，经济会不会自己回到充分就业。在弗里德曼的论证里，长期会回到自然率，短期里名义工资按预期的物价定，跟不上实际的物价（本书的概括）。\chref{ch:islm}的 IS–LM 把货币工资当作给定的，附加预期的曲线让它随预期的通胀调整（本书的解读）。

\begin{pitfall}[自然失业率就是应该有的失业率]
弗里德曼说，“自然”不是说它一成不变，决定它的许多市场特征是人为的、是政策造成的\cite{friedman1968}。菲尔普斯说，在他的模型里，均衡时也有非自愿的失业，需求的意外增加会减轻这种沉重的失业\cite{phelps2006}。自然率只是货币政策在长期里改变不了的那个失业率，高低取决于劳动市场的制度，不一定是社会想要的（本书的解读）。制度怎样影响失业，在\chref{ch:labor}。
% 据: f68_immutable, f68_manmade, ph_involuntary, ph_moderates
\end{pitfall}

\section{附加预期的菲利普斯曲线}\label{sec:naturalrate-model}

把这个论证写成一个式子。本书用$\pi_t$记第$t$年的通胀率，$\pi^{e}_t$记人们对这一年的通胀的预期，$u_t$记失业率，$u^{n}$记自然失业率，$\alpha>0$记失业对通胀的影响。弗里德曼 1976 年把自然失业率记作$U_N$\cite{friedman1976}。附加预期的菲利普斯曲线是
\begin{keyeq}
\pi_t=\pi^{e}_t-\alpha\,(u_t-u^{n}) .\label{eq:naturalrate-curve}
\end{keyeq}
预期的通胀不变，它就是\chref{ch:phillips}的菜单，失业越低，通胀越高。预期一变，整条曲线跟着上下平移，每一种预期对应一条短期曲线（本书的推导）。菜单漂移，就是曲线随预期移动。设预期是 2\%，自然率是 6\%（教学用的假设），失业 3\% 时通胀是$2+3\alpha$。人们照这个通胀预期下一年，同一个 3\% 的失业配的通胀就多出$3\alpha$（本书的推导）。弗里德曼 1976 年说，这个假说也叫自然率假说或加速论，已被经济学家广泛接受，但远不是人人都接受\cite{friedman1976}。
% 据: f76_un, f76_names, f76_accepted

预期为什么要紧？弗里德曼 1976 年把理由讲得更直白。雇用多是长期的约定，价格与数量都只能逐步调整，订约时看的不只是眼下的价格，也是约定期内预期的价格\cite{friedman1976}。他说，只有意外才起作用\cite{friedman1976}。如果人人都预期物价每年涨 20\%，这个预期就会写进以后的工资合同，实际工资的变动与人人都预期物价不涨时完全一样\cite{friedman1976}。20\% 的通胀也就没有理由对应另一个失业水平\cite{friedman1976}。
% 据: f76_commitments, f76_commitments2, f76_surprises, f76_twenty, f76_twenty2

人们怎样形成预期？最简单的假设是适应性预期，即按去年的通胀预期今年，$\pi^{e}_t=\pi_{t-1}$。代入式~\eqref{eq:naturalrate-curve}，得
\begin{keyeq}
\pi_t-\pi_{t-1}=-\alpha\,(u_t-u^{n}) .\label{eq:naturalrate-accel}
\end{keyeq}
失业低于自然率，通胀逐年上升。失业高于自然率，通胀逐年下降。只有失业等于自然率，通胀才不变（本书的推导）。这就是加速论。长期里通胀不再变动，失业只能等于$u^{n}$，通胀停在哪个水平都一样，长期的曲线是$u^{n}$处的一条竖线（本书的推导）。人们若会预期政策本身，结论还要更进一步，这是\chref{ch:ratex}的理性预期。

弗里德曼 1976 年用一张图讲同一个过程\cite{friedman1976}。从稳定的 E 点出发，通胀从 A 升到 B 并停在那里\cite{friedman1976}。失业先沿着预期通胀为 A 的短期曲线降到 F 点\cite{friedman1976}。预期调整以后，短期曲线逐渐上移到预期通胀为 B 的那一条，失业也逐渐从 F 移到 G，回到原来的水平\cite{friedman1976}。他承认这个分析过于简化，只假定了一次意外，没有处理时滞、超调与预期的形成\cite{friedman1976}。
% 据: f76_fig2_start, f76_fig2_f, f76_fig2_g, f76_back, f76_oversimplified

预期为什么会慢慢变？弗里德曼 1976 年讲过英美的经验\cite{friedman1976}。二战以前，英国两百年、美国一个半世纪里，物价围绕一个大体不变的水平波动，战时大涨，战后跌回战前的水平\cite{friedman1976}。“正常”的物价水平深深嵌在两国的金融制度与公民的习惯里\cite{friedman1976}。二战刚结束时，战时的通胀之上又叠加了战后的通胀，两国人预期的却是通缩\cite{friedman1976}。他说，对战后通缩的担心过了很久才消散，预期向货币制度的根本变化调整，至今还远没有完成\cite{friedman1976}。
% 据: f76_centuries, f76_normal, f76_deflation, f76_dissipate, f76_incomplete

弗里德曼 1976 年列举了经济学家当时的几种态度\cite{friedman1976}。少数人仍守着原来的菲利普斯曲线\cite{friedman1976}。更多的人承认短期与长期的曲线不同，却认为长期的曲线也向下倾斜，只是更陡\cite{friedman1976}。还有人用通胀的加速与失业之间的稳定关系，代替通胀与失业之间的稳定关系，式~\eqref{eq:naturalrate-accel}就是这一种\cite{friedman1976}。他提醒，把人们推到二阶导数的那套逻辑，也可能把人们推向更高阶的导数\cite{friedman1976}。
% 据: f76_cling, f76_steeper, f76_acceleration, f76_derivatives

\begin{example}{守住 3\% 的失业}\label{ex:naturalrate-accel}
设$\alpha=0.5$，$u^{n}=6\%$（教学用的假设），通胀原来稳定在 2\%。货币当局从第 1 年起把失业守在弗里德曼举例的 3\%。（1）按式~\eqref{eq:naturalrate-accel}，前四年的通胀各是多少？（2）若人们的预期一直停在 2\%，结果怎样？（3）第 4 年以后放手，失业回到 6\%，通胀停在哪里？
\solution
（1）每年上升$0.5\times(6-3)=1.5$个百分点，依次是 3.5\%、5\%、6.5\% 与 8\%。（2）按式~\eqref{eq:naturalrate-curve}，$\pi^{e}_t$不变，每年都是$2+1.5=3.5\%$，这就是菜单上的一个点。（3）失业等于自然率，通胀不再变动，停在 8\%（本书的计算）。失业回到了原处，留下的是更高的通胀。
\end{example}
% 据: f68_target

\begin{compare}{短期菲利普斯曲线与长期菲利普斯曲线}{短期曲线}{长期曲线}
\cmprow{预期}{给定，$\pi^{e}_t$不变}{追上了实际，$\pi^{e}=\pi$}
\cmprow{形状}{向下倾斜，斜率是$-\alpha$}{$u^{n}$处的一条竖线}
\cmprow{取舍}{总有暂时的取舍\cite{friedman1968}}{没有永久的取舍\cite{friedman1968}}
\cmprow{多久}{最初的效果持续两到五年\cite{friedman1968}}{完全调整大约要二十年\cite{friedman1968}}
\cmprow{政策能做什么}{上升的通胀可以降低失业\cite{friedman1968}}{高通胀不能\cite{friedman1968}}
\cmpsum{两条曲线是同一个式~\eqref{eq:naturalrate-curve}，预期给定时看到短期的曲线，预期追上来以后看到长期的曲线（本书的概括）。}
% 据: f68_tradeoff, f68_rising, f68_years, f68_decades
\end{compare}

\section{反通胀的代价}\label{sec:naturalrate-cost}

反过来，要让通胀降下来，就得让失业高于自然率。把式~\eqref{eq:naturalrate-accel}从第 1 年加到第$T$年，左边逐项相消，得
\begin{keyeq}
\sum_{t=1}^{T}(u_t-u^{n})=\frac{\pi_0-\pi_T}{\alpha} .\label{eq:naturalrate-sacrifice}
\end{keyeq}
左边是累计多出的失业，单位是百分点·年。通胀要降多少，累计就要多失业多少，与降得快慢无关（本书的推导）。换成产出，要用\chref{ch:phillips}的奥肯定律，失业率每高出 1 个百分点，产出约少 3.2\%（式~\eqref{eq:phillips-okun}）\cite{okun1962}。下面把奥肯式子里的 4\% 换成$u^{n}$，这是教学用的假设。通胀每降 1 个百分点，累计要少生产一年产出的$3.2/\alpha$个百分点，这个比率叫作牺牲率（本书的推导）。

\begin{example}{慢降与急降}\label{ex:naturalrate-disinflation}
仍设$\alpha=0.5$，$u^{n}=6\%$。通胀原来稳定在 10\%，要降到 4\%。（1）渐进：失业守在 8\%，要几年？每年的通胀是多少？（2）急降：失业守在 12\%，要几年？（3）两种办法累计多出的失业各是多少？按奥肯的系数 3.2，累计少生产多少？牺牲率是多少？
\solution
（1）每年降$0.5\times(8-6)=1$个百分点，要 6 年，依次是 9\%、8\%、7\%、6\%、5\% 与 4\%。（2）每年降$0.5\times(12-6)=3$个百分点，要 2 年，依次是 7\% 与 4\%。（3）都是$(10-4)/0.5=12$个百分点·年，与式~\eqref{eq:naturalrate-sacrifice}相合。累计少生产$3.2\times12=38.4\%$个年产出，牺牲率是$38.4/6=6.4$（本书的计算）。快慢只改变代价在各年的分布，总量不变。
\end{example}

\fig{naturalrate_sim}{慢降与急降（\exref{ex:naturalrate-disinflation}，本书的计算）。$\alpha=0.5$，$u^{n}=6\%$，通胀从 10\% 降到 4\%，预期按上一年的通胀形成。(a) 失业率（\%），虚线是自然率。(b) 通胀率（实线）与预期的通胀率（虚线），\%。两种办法累计多出的失业都是 12 个百分点·年。}{fig:naturalrate_sim}

教学的数只说明机制，实际的斜率要从数据里估计。1978 年，奥肯（Arthur Okun，1928—1980）综合了多种计量的证据，得出一个悲观的结论\cite{gordon1982,loc44,gnd44}。按 Gordon 与 King 的转述，他认为通胀每永久降低 1 个百分点，累计要牺牲一年 GNP 的 10\%\cite{gordon1982}。他算出的代价在一年 GNP 的 6\% 到 18\% 之间，平均是 10\%\cite{gordon1982}。按上面的模型，牺牲率是 10，$\alpha$约为 0.32（本书的计算）。
% 据: gk_okun, gk_okun_range, gk_fn1, loc_okun, gnd_okun

Gordon 与 King 指出，奥肯用的是反通胀实验头一年的产出损失与通胀降幅之比，没有考虑这个比率在实验进行中的变化\cite{gordon1982}。他们自己模拟了 1981—1992 年的两条路径\cite{gordon1982}。对照的路径让 M1 每年增长 6.6\%\cite{gordon1982}。减速的路径让它 1981 年增长 5.0\%，1982 年增长 4.0\%\cite{gordon1982}。此后每年减 0.5 个百分点，1986—1992 年保持 2.0\%\cite{gordon1982}。他们说，减速的路径近似联邦储备委员会当时的官方政策\cite{gordon1982}。两条路径的产出之差是代价，通胀之差是收益（本书的解读）。
% 据: gk_firstyear, gk_simulate, gk_control, gk_volcker, gk_volcker2, gk_official

按他们的基本模型，累计的产出损失是一年 GNP 的 13.4\%，通胀平均每年降 4.4 个百分点\cite{gordon1982}。不贴现的牺牲率是 3.0，贴现以后是 4.3\cite{gordon1982}。他们承认，只算到 1992 年，漏掉了此后通胀较低的好处，会高估牺牲率\cite{gordon1982}。也漏掉了超调以后的衰退、产出与通胀的不稳定，以及 1981—1985 年萧条中投资减少的损失，会低估它\cite{gordon1982}。他们假定两方面的净影响很小，不会让结果明显偏离\cite{gordon1982}。
% 据: gk_basic, gk_ratios, gk_over, gk_under, gk_net

代价落在失业的人身上。按表 B-40，1975 年全体劳动者的失业率是 8.5\%，黑人是 14.8\%\cite{erp2000}。16—19 岁的青年是 19.9\%\cite{erp2000}。1982 年全体升到 9.7\%，黑人是 18.9\%，与妻子同住的已婚男子只有 6.5\%\cite{erp2000}。青年是 23.2\%\cite{erp2000}。从 1979 年到 1982 年，全体的失业率升了 3.9 个百分点，黑人升了 6.6 个（本书的计算）。青年升了 7.1 个（本书的计算）。决定紧缩的是货币当局，失业升得最多的是本来就最容易失业的人（本书的概括）。沃尔克 1979—1982 年的反通胀实际付出了多少，在\chref{ch:credibility}。
% 据: b40_1975, b40_1982, b40_1979, b40_cols

\section{1970 年代的检验}\label{sec:naturalrate-test}

先看美国。\figref{fig:naturalrate_us}把 1960—1985 年美国的失业率与消费物价的年涨幅画在一起\cite{erp2000}。1960—1969 年的点像一张菜单，失业低的年份通胀高，相关系数约为$-0.84$（本书的计算）。1970—1985 年的点散开了，相关系数只有约$-0.11$（本书的计算）。这十六年平均失业 6.9\%，通胀 7.0\%（本书的计算）。前十年只有 4.8\% 与 2.4\%（本书的计算）。
% 据: b40_title, b62_title

顺着年份看，1970 年代的点在图上绕圈（本书的解读）。1974 年通胀跳到 11.0\%，失业还只有 5.6\%\cite{erp2000}。1975 年失业升到 8.5\%，次年通胀降到 5.8\%\cite{erp2000}。1980 年通胀升到 13.5\%，失业是 7.1\%\cite{erp2000}。1982 年失业升到 9.7\%，次年通胀降到 3.2\%\cite{erp2000}。后一圈比前一圈更靠右，同样的通胀配着更高的失业（本书的解读）。
% 据: b40_1974, b62_1974, b40_1975, b62_1976, b40_1980, b62_1980, b40_1982, b62_1983

\fig{naturalrate_us}{美国的失业率与通胀率，1960—1985 年。横轴是全体平民劳动者的失业率（\%），纵轴是全体城市消费者的消费物价指数的年涨幅（\%），相邻年份连线。蓝点是 1960—1969 年，橙点是 1970—1985 年。数据：《总统经济报告》（2000）表 B-40 与 B-62，原始来源是美国劳工统计局。}{fig:naturalrate_us}

加速论说的是通胀的变动。把每年通胀的变动对失业率回归，1960—1985 年得$\pi_t-\pi_{t-1}=4.64-0.74\,u_t$（本书的计算）。通胀不变时失业率约为 6.3\%，这是按式~\eqref{eq:naturalrate-accel}估出的自然率（本书的计算）。失业低于 5\% 的 7 年里，即 1965—1970 年与 1973 年，通胀年年上升（本书的核对）。通胀下降的 9 年，失业都在 5.6\% 以上（本书的核对）。可是失业在 7\% 以上的 9 年里，1977、1980 与 1984 年的通胀仍在上升（本书的核对）。通胀的变动不只取决于失业，油价的冲击是\chref{ch:oil}的题目。

弗里德曼在演讲里用七个工业国的数据作了同样的检验（\tabref{tab:naturalrate-seven}）\cite{friedman1976}。按五年平均，前两个五年之间，七国中有五国的通胀与失业反向变动，这是简单的菲利普斯曲线预言的结果\cite{friedman1976}。第二与第三个五年之间只剩四国\cite{friedman1976}。最后两个五年之间只剩一国，就是意大利\cite{friedman1976}。他说意大利并不是真正的例外，1971—1975 年它的失业比前五年略低，通胀却增加到原来的三倍多，1973 年以后两者都急剧上升\cite{friedman1976}。
% 据: f76_table, f76_five, f76_four, f76_one, f76_italy

\begin{table}[tbp]
\centering\small
\caption{七国的通胀与失业，1956—1975 年（五年平均，\%）}\label{tab:naturalrate-seven}
\begin{tabularx}{\linewidth}{@{}Y rr rr rr rr@{}}
\toprule
 & \multicolumn{2}{c}{1956—1960} & \multicolumn{2}{c}{1961—1965} & \multicolumn{2}{c}{1966—1970} & \multicolumn{2}{c@{}}{1971—1975} \\
 & 通胀 & 失业 & 通胀 & 失业 & 通胀 & 失业 & 通胀 & 失业 \\
\midrule
法国 & 5.6 & 1.1 & 3.7 & 1.2 & 4.4 & 1.7 & 8.8 & 2.5 \\
德国 & 1.8 & 2.9 & 2.8 & 0.7 & 2.4 & 1.2 & 6.1 & 2.1 \\
意大利 & 1.9 & 6.7 & 4.9 & 3.1 & 3.0 & 3.5 & 11.3 & 3.3 \\
日本 & 1.9 & 1.4 & 6.2 & 0.9 & 5.4 & 1.1 & 11.4 & 1.4 \\
瑞典 & 3.7 & 1.9 & 3.6 & 1.2 & 4.6 & 1.6 & 7.9 & 1.8 \\
英国 & 2.6 & 1.5 & 3.5 & 1.6 & 4.6 & 2.1 & 13.0 & 3.2 \\
美国 & 2.0 & 5.2 & 1.3 & 5.5 & 4.2 & 3.9 & 6.7 & 6.1 \\
七国平均 & 2.8 & 3.0 & 3.7 & 2.0 & 4.1 & 2.2 & 9.3 & 2.9 \\
\bottomrule
\end{tabularx}
\datasource{弗里德曼（1976）表 1。通胀是消费物价的年复合涨幅，按 1955—1960 年、1960—1965 年等五年计算；失业是五年的平均，所以通胀的时点比失业早半年。七国平均等于七国的简单平均（本书的核对）。}
\end{table}
% 据: f76_t1956, f76_t1961, f76_t1966, f76_t1971, f76_t1note

他也承认，这些数据至多只有提示的作用\cite{friedman1976}。七国受共同的国际因素影响，不是七组独立的证据\cite{friedman1976}。石油危机同时冲击了七国，可他说，五年的平均与逐年的数据都表明，光凭它解释不了这个现象\cite{friedman1976}。日本的失业比通胀升得少得多，他认为这反映了在日本的制度环境里，失业的含义不同\cite{friedman1976}。
% 据: f76_suggestive, f76_common, f76_oil, f76_oil2, f76_japan

逐年的数据讲的是同一个故事，只是更乱\cite{friedman1976}。早些年各国差别很大，意大利的物价与失业几乎没有关系，英国与美国逐年的反向关系相当清楚\cite{friedman1976}。近几年，法国、美国、英国、德国与日本的通胀与失业都明显上升，只有瑞典与意大利不合这个模式\cite{friedman1976}。他还猜测，抬高美国自然率的那些力量，也可能在别的国家起作用\cite{friedman1976}。
% 据: f76_annual, f76_early, f76_recent, f76_sweden, f76_forces

通胀高了，失业为什么也高？弗里德曼说，近年来，在几年的时段里，更高的通胀常常伴随着更高而不是更低的失业，简单的统计曲线似乎向上倾斜\cite{friedman1976}。他的初步假说是，向下倾斜的短期曲线会随预期的调整消失，向上倾斜的曲线也可能是一种过渡现象，要等人们不只调整预期，也调整制度与政治安排才会消失\cite{friedman1976}。
% 据: f76_stage3, f76_transitional

在弗里德曼的假说里，竖直的长期曲线对应的是完全预期到的通胀\cite{friedman1976}。预期的通胀若是每年 20\%，长期工资合同会让工资每年比零通胀时多涨 20\%，长期贷款的利率高出 20 个百分点，等于把所有合同都指数化了\cite{friedman1976}。这要求通胀稳定，价格能自由调整，合同的指数化没有障碍\cite{friedman1976}。他说，20\% 的通胀若持续几十年，这些条件可以大体满足，所以他仍保留长长期的竖直曲线\cite{friedman1976}。可是一国刚转向更高的通胀时，这些条件会系统地被违背，这样的过渡期可能长达几十年\cite{friedman1976}。
% 据: f76_anticipated, f76_index20, f76_conditions, f76_conditions2, f76_longlong, f76_transition

他说，通胀不只更高，也越来越不稳定，政府对定价的干预也越来越多\cite{friedman1976}。通胀越不稳定，就越难从绝对价格里读出相对价格的信号\cite{friedman1976}。这些都会降低经济的效率，很可能抬高记录的失业率\cite{friedman1976}。
% 据: f76_volatile, f76_signal, f76_frictions

通胀的代价也不是人人一样。弗里德曼说，在制度与金融合约都按长期“正常”物价水平安排的国家，通胀让一些人得益，如房主，让另一些人受损，如储蓄存款与固定利息证券的持有人\cite{friedman1976}。社会因此分化，政治不安增加，政府治理的能力下降，要求强力行动的压力却在加大\cite{friedman1976}。他判断这种局面不能持续，要么恶化成恶性通胀，要么制度适应长期的通胀，要么政府转向低通胀、少干预定价的政策\cite{friedman1976}。
% 据: f76_groups, f76_polarized, f76_cannotlast

莫迪利亚尼与 Papademos 给了这个失业率另一个名字。他们把它叫作“不引起通胀的失业率”，简称 NIRU，定义是只要失业高于它，通胀就可望下降，起点很低时也许例外\cite{modigliani1975}。他们说，主张曲线垂直与不垂直的两派，都隐含着 NIRU 的存在\cite{modigliani1975}。按他们的说法，1953—1971 年的数据与一个界定清楚、长期稳定的 NIRU 相符，前提是失业的量度要按劳动力构成的变化调整\cite{modigliani1975}。
% 据: mp_niru, mp_schools, mp_stable

他们的证据是一张图\cite{modigliani1975}。横轴是经济顾问委员会按 1956 年的劳动力构成调整过的失业率，纵轴是那一年通胀率的变化，通胀用不含食品的消费物价指数来量\cite{modigliani1975}。图上的纵轴原先标错了，另附的一页勘误作了更正\cite{modigliani1975}。调整后的失业率高于 5\% 的 8 年里，除了 1962 年，通胀都下降了，他们推断这个失业率的 NIRU 至多是 5\%\cite{modigliani1975}。等于或低于 4.3\% 的 9 年里，除了两年，通胀都明显上升\cite{modigliani1975}。他们用回归估出 4.88\%\cite{modigliani1975}。按当时的劳动力构成，还要加 0.8 个百分点\cite{modigliani1975}。按官方的失业率，NIRU 当时估计略高于 5.5\%\cite{modigliani1975}。
% 据: mp_1956, mp_fig1, mp_pcx, mp_erratum, mp_five, mp_43, mp_488, mp_08, mp_55

他们也提醒，有证据表明，过去二十年 NIRU 被贸易条件的有利趋势压低了，即私人非农部门的价格相对于进口品与农产品的价格\cite{modigliani1975}。这个趋势一旦终止或逆转，NIRU 至少会暂时上升\cite{modigliani1975}。
% 据: mp_terms, mp_reversal

他们拿 NIRU 给下一年的货币政策定目标\cite{modigliani1975}。他们建议，把 1977 年年中的失业目标稳妥地定在 6\%，为此产出在今后两年至少要增长 17\%\cite{modigliani1975}。头一年货币收入的年增长要超过 15\%\cite{modigliani1975}。即使复苏主要靠财政政策推动，货币供给也要有一段时间增长 10\% 以上\cite{modigliani1975}。人们担心货币这样加速会推高通胀\cite{modigliani1975}。他们回答，证据表明货币只通过总需求与就业间接影响通胀\cite{modigliani1975}。守着 1974 年那样的货币增长目标，复苏很可能乏力，失业上升，甚至再次衰退\cite{modigliani1975}。货币数量与通胀的关系，是\chref{ch:monetarism}的题目。
% 据: mp_target, mp_17, mp_income15, mp_money10, mp_concern, mp_indirect, mp_sluggish

\begin{compare}{自然失业率与 NIRU}{自然失业率（弗里德曼，1968）}{NIRU（莫迪利亚尼与 Papademos，1975）}
\cmprow{定义}{瓦尔拉斯一般均衡方程组算出的失业水平\cite{friedman1968}}{失业高于它，通胀就可望下降\cite{modigliani1975}}
\cmprow{长期的曲线}{垂直，没有永久的取舍\cite{friedman1968}}{垂直与不垂直的两派都隐含它\cite{modigliani1975}}
\cmprow{数值}{货币当局不知道它是多少\cite{friedman1968}}{按官方失业率略高于 5.5\%\cite{modigliani1975}}
\cmprow{会不会变}{不是一成不变的，许多决定因素是人为的\cite{friedman1968}}{1953—1971 年稳定，前提是按劳动力构成调整\cite{modigliani1975}}
\cmprow{用来做什么}{说明货币政策钉不住失业率\cite{friedman1968}}{给下一年的货币政策定失业与产出的目标\cite{modigliani1975}}
\cmpsum{两个名字指的是同一个位置。自然率讲的是实际工资结构的均衡，NIRU 只讲失业高于它时通胀会降（本书的概括）。}
% 据: f68_natural, mp_niru, f68_tradeoff, mp_schools, f68_cannot_know, mp_55, f68_manmade, mp_stable, mp_target
\end{compare}

英国的数据更极端（\figref{fig:naturalrate_jpuk}a）。按英格兰银行千年数据集里的统计，1955—1969 年英国的失业率在 1.6\% 到 3.5\% 之间\cite{boe2017}。消费物价每年涨 0.6\% 到 5.2\%\cite{boe2017}。1971 年物价涨了 9.4\%，失业是 4.1\%\cite{boe2017}。1974 年物价涨了 15.7\%，失业只有 3.6\%\cite{boe2017}。1975 年物价涨了 22.7\%，失业是 4.5\%\cite{boe2017}。1980 年以后失业急升，1981 年是 9.6\%，1983 年到了 11.5\%\cite{boe2017}。物价的涨幅 1980 年仍有 15.2\%，1983 年降到 4.9\%\cite{boe2017}。卡拉汉的那番话，说在两者之间（本书的解读）。

\fig{naturalrate_jpuk}{英国与日本的失业率与通胀率（\%）。(a) 英国，1955—1985 年：横轴是失业人数占就业与失业人数之和的百分比，纵轴是消费物价的年涨幅，数据据英格兰银行《英国千年宏观经济数据》（3.1 版）表 A50 与 A47。表中的就业与失业，1970 年以前据《经济趋势年度增刊》（1997 年版）的年中估计，女性失业按劳动力调查与申领人数之差调整过，1971 年起是英国国家统计局劳动力调查的 MGRZ 与 MGSC 序列。消费物价 1949 年以后是英国国家统计局的 CPI。(b) 日本，1955—1985 年：横轴是完全失业率，纵轴是不含自有住房估算房租的消费物价指数的年涨幅；出典：「労働力調査」（総務省統計局）、「消費者物価指数」（総務省統計局）を加工して作成。1973 年以前的失业率不含冲绳，1973 年取含冲绳的数。}{fig:naturalrate_jpuk}

日本的点挤在图的左边（\figref{fig:naturalrate_jpuk}b）。按劳动力调查，1955—1972 年日本的完全失业率在 1.1\% 到 2.5\% 之间\cite{estatlfs}。同一时期，消费物价的年涨幅最低是负的，最高到 7.7\%\cite{estatcpi}。失业率最高与最低只差 1.4 个百分点（本书的计算）。1974 年物价飞涨，失业却几乎没动，下面的案例细看这件事。

\begin{case}{日本 · 1973—1974 · 物价飞涨，失业不动}
按总务省统计局的劳动力调查，1960—1973 年日本的完全失业率在 1.1\% 到 1.7\% 之间，1974 年是 1.4\%\cite{estatlfs}。物价却飞涨。不含自有住房估算房租的消费物价，1973 年涨 11.7\%，1974 年涨 24.5\%\cite{estatcpi}。含这一项的综合指数，1974 年涨 23.2\%\cite{estatcpi}。
% 据: lfs_table, cpi_xr_table, cpi_all_table

经济企划厅在 1974 年 8 月的年度经济报告里分析了原因\cite{epa1974}。它说，批发物价从 1972 年 8 月起急涨，消费物价从 1973 年 3 月起急涨\cite{epa1974}。按它的分解，1973 年度批发物价涨 35.4\%，约四成来自国内供需的紧张，约六成来自原油等初级产品价格暴涨这些海外因素\cite{epa1974}。工资成本的影响“极小”，“通胀心理”却在加强\cite{epa1974}。
% 据: wp_date, wp_wholesale, wp_consumer, wp_354, wp_wages, wp_psychology

它说，最根本的原因是对日本经济成长轨道改变的适应太迟\cite{epa1974}。它把昭和 40 年代（1965 年开始的十年）分成前后两半来比，后一半民间设备投资的实际增长从 16.1\% 降到 10.4\%\cite{epa1974}。制造业生产能力的增长从 12.1\% 降到 7.4\%\cite{epa1974}。名义的需求却从 16.5\% 加快到 17.4\%\cite{epa1974}。货币供给（M2）的增长从 16.6\% 升到 21.9\%，1971—1972 年的汇率调整期增势更强，部分是为了让经济尽快走出美元冲击带来的萧条\cite{epa1974}。按它的说法，问题已经从需求不足变成供给能力不足\cite{epa1974}。
% 据: wp_basic, wp_halves, wp_invest, wp_capacity, wp_demand, wp_m2, wp_m2b, wp_dollar, wp_supply

决定压住需求的是日本银行与政府（本书的概括）。企划厅说，1973 年度的政策以总需要抑制为轴，比过去强力得多\cite{epa1974}。公定贴现率从 4 月到 12 月提高 5 次，从 4.25\% 升到 9\%\cite{epa1974}。过去紧缩时，提高的幅度只在 1\% 左右\cite{epa1974}。日本银行的窗口指导更严，对象扩大到相互银行与较大的信用金库\cite{epa1974}。存款准备率也在一年里提高了 5 次，公共投资逐步推迟\cite{epa1974}。
% 据: wp_axis, wp_discount, wp_past, wp_window, wp_reserve, wp_public

付代价的是收入低的家庭。企划厅说，物价飞涨对收入低的阶层影响更大\cite{epa1974}。收入高的阶层实际支出继续增加，收入低的阶层进入 1974 年以后，连生活必需的实际支出也减少了\cite{epa1974}。政府给手纸、合成洗涤剂等抢购的商品定了标准价格，10 月又决定把生活扶助的标准提高 5\%\cite{epa1974}。
% 据: wp_lowincome, wp_classes, wp_tissue, wp_assistance

自然率的读法能说明一件事，日本的失业很低，物价飞涨并没有换来更低的失业（本书的解读）。它说明不了失业为什么几乎不动，弗里德曼也只说日本的失业含义不同\cite{friedman1976}。1975—1978 年，通胀从 11.8\% 降到 3.8\%\cite{estatcpi}。失业只从 1.9\% 升到 2.2\%\cite{estatlfs}。按失业率量，这次反通胀的代价很小，按产出或工时量是多少，本章的数据回答不了（本书的解读）。
% 据: f76_japan, lfs_table, cpi_xr_table
\end{case}

\section{争论}\label{sec:naturalrate-debate}

\subsubsection{是不是同一个发现}

菲尔普斯（Edmund S. Phelps，1933—2026）从 1960 年代中期起研究失业怎样决定\cite{phelps2006,loc44,gnd44}。他在 2006 年的诺贝尔奖演讲里说，起初推动他的问题是，有效需求的激增为什么会增加产出与就业，为什么不只是物价与货币工资跳上去\cite{phelps2006}。他接着问，均衡里怎么会有正的非自愿失业\cite{phelps2006}。按他的模型，要是没有失业，雇员辞职会多到每家企业都想用更高的工资压过别家，好省下人员流动带来的大笔培训开支\cite{phelps2006}。他说，这个预期的框架第一次发表的应用，是他 1967 年为最优的通胀政策建的模型\cite{phelps2006}。这篇论文 1966 年初写于伦敦经济学院，针对的是当时用菲利普斯曲线求最优通胀率的做法\cite{phelps2006}。弗里德曼 1976 年说，区分意外的需求变动的短期效果与长期效果的这个假说，特别是由菲尔普斯与他自己提出的\cite{friedman1976}。
% 据: ph_mid1960s, ph_question, ph_involuntary, ph_quitting, ph_1967, ph_lse, ph_okun1965, f76_phelps, loc_phelps, gnd_phelps

两人讲的误判也不一样（本书的解读）。弗里德曼 1976 年的说法是，需求意外加速，每个生产者都把它至少部分当作只对自己产品的需求，愿意付更高的名义工资\cite{friedman1976}。工人在意的是工资对一般物价的购买力，会把名义工资的上涨看成实际工资的上涨，于是多供给劳动\cite{friedman1976}。在菲尔普斯 1968 年的模型里，每家企业都误以为需求的增加多半只是它自己的，于是低估了别家工资与价格的上涨，自己的工资与价格只提一点\cite{phelps2006}。雇员以为别家的工资比本家低，辞职减少，就业随之扩大\cite{phelps2006}。
% 据: f76_producer, f76_workers, f76_workers2, ph_1968, ph_unique, ph_little, ph_quit_less

菲尔普斯不同意把两人并列。他说，他的模型常被看作与弗里德曼 1968 年的模型基本相同，两者被当作同一件事的同时发现\cite{phelps2006}。他认为两人发现的是“两种不同的现象”。弗里德曼的模型讲的是劳动参与的自然率，他的讲的是失业的自然率\cite{phelps2006}。在弗里德曼的模型里，没有被察觉的需求增加是对竞争均衡的一种不受欢迎的偏离。在他的模型里，它减轻了本来很沉重的非自愿失业\cite{phelps2006}。这是同一个发现还是两个，至今没有定论。
% 据: ph_identical, ph_distinct, ph_participation, ph_unwelcome, ph_moderates

\subsubsection{自然率会不会变}

弗里德曼说，自然率不是一个数值常数\cite{friedman1976}。他举了美国的自然率明显上升的两个原因。一是妇女、青少年与非全日工在劳动力中的比重越来越大，他们进出劳动市场、换工作都更频繁\cite{friedman1976}。二是失业保险与其他的失业救济覆盖了更多的工人\cite{friedman1976}。莫迪利亚尼与 Papademos 则说，按劳动力构成调整以后，NIRU 在 1953—1971 年是稳定的\cite{modigliani1975}。Gordon 检验了按构成加权的自然失业率在 1970 年代上移的假说，结果是拒绝\cite{gordon1982}。
% 据: f76_constant, f76_two, f76_women, f76_insurance, mp_stable, gk_fn20

失业的高低也不等于效率的高低。弗里德曼说，把记录的高失业当作资源用得没效率的证据，是严重的错误\cite{friedman1976}。低失业可能说明经济被强行驱动，工人误以为实际工资会更高，放弃了本来更看重的闲暇\cite{friedman1976}。僵硬的经济可能给每人一个固定的位置，充满变化的经济自然率反而可能高\cite{friedman1976}。他举例说，大约 1950—1970 年，日本与英国的平均失业都很低，日本却增长迅速，英国停滞\cite{friedman1976}。用本章的回归，1960—1969 年估出的自然率约为 5.8\%，1970—1985 年约为 6.8\%（本书的计算）。弗里德曼自己说，货币当局不知道自然率是多少\cite{friedman1968}。它是常数、慢慢移动的数，还是根本估不准，至今没有定论。
% 据: f76_error, f76_forced, f76_rigid, f76_japan_uk, f68_cannot_know

\subsubsection{反通胀要付多大代价}

奥肯的估计在一年 GNP 的 6\% 到 18\% 之间，平均 10\%\cite{gordon1982}。Gordon 与 King 1982 年把牺牲率定义为永久消除 1 个百分点的通胀所要损失的产出，结论是只有奥肯的约一半\cite{gordon1982}。他们的基本模型给出的贴现牺牲率是 4.3\cite{gordon1982}。把国际变量逐个当作外生，它依次升到 5.8、7.2 与 9.9\cite{gordon1982}。最后一个已接近奥肯的平均数\cite{gordon1982}。他们偏爱 5.8 这个数\cite{gordon1982}。按它，通胀长期降低 5 个百分点，要放弃现值约为一年自然 GNP 29\% 的产出，按当时的价格将近 1 万亿美元\cite{gordon1982}。
% 据: gk_okun_range, gk_half, gk_43, gk_99, gk_prefer, gk_trillion

产出的损失还要与通胀降低的好处对比\cite{gordon1982}。他们说，牺牲率若是 6，代价就相当于一年自然 GNP 的 30\%，平均每人约 4000 美元\cite{gordon1982}。算上失业者多出的闲暇，1 万亿美元的产出损失减为约 8600 亿美元\cite{gordon1982}。好处方面，Fischer 估计通胀降低 5 个百分点，每年的收益约为 GNP 的 0.30\%\cite{gordon1982}。把通胀对储蓄的影响算进去，收益可以高到 GNP 的 1\%，可这要假定税制改不了\cite{gordon1982}。他们认为，除非假定不中性的税收扭曲与金融管制永远存在，贴现的福利收益不大可能接近 1 万亿美元，而这个假定不可信\cite{gordon1982}。
% 据: gk_compare, gk_six, gk_leisure, gk_fischer, gk_saving, gk_welfare

他们说，美国政府当时的官方货币政策给全世界造成的产出损失，比这 1 万亿美元还多\cite{gordon1982}。他们也认为，公众厌恶通胀，在很大程度上是把通胀本身的影响与 1970 年代油价冲击、生产率放慢造成的实际收入损失混在了一起\cite{gordon1982}。参加讨论的 Sims 指出文章没有考虑财政政策\cite{gordon1982}。他认为，联邦储备坚决的反通胀配上历史上很高而且还在增长的赤字，这一政策组合的牺牲率可能比他们的估计高得多\cite{gordon1982}。反通胀的代价究竟多大，至今没有定论。
% 据: gk_world, gk_aversion, gk_sims, gk_sims2

\begin{deeper}[史料说明]
\textbf{版本。}弗里德曼 1968 年的演说据美国经济学会网站的 PDF，是原刊的扫描本，文字层把一些词连在一起或读错。弗里德曼与菲尔普斯的诺贝尔奖演讲据诺贝尔奖网站的 PDF。莫迪利亚尼与 Papademos（1975）、Gordon 与 King（1982）据布鲁金斯学会网站上的扫描本。美国的失业率与消费物价据 GovInfo 上 2000 年《总统经济报告》的表 B-40 与 B-62。日本的数据据 e-Stat 上劳动力调查的长期时间序列表与 2015 年基准消费物价指数的长期时间序列。经济企划厅 1974 年的报告据内阁府网站的网页版。

\textbf{转述与间接的证词。}奥肯（1978）没有读，正文只按 Gordon 与 King 的转述。菲尔普斯 1967 年与 1968 年的论文没有读，只按他 2006 年的演讲与弗里德曼 1976 年的演讲。卡拉汉的讲话只据弗里德曼的引文。

\textbf{没有取得的。}奥肯（1978）在圣路易斯联邦储备银行的 FRASER 上，2026 年 11 月 12 日那里没有应答。莫迪利亚尼 1944 年的论文仍然没有取得（\chref{ch:islm}）。后来通行的 NAIRU 一词怎样从 NIRU 演变而来，本书没有找到可靠的第一手说明，正文只用 NIRU。

\textbf{几个证人不合。}1969 年美国的消费物价，\chref{ch:phillips}按城市工薪劳动者的消费物价指数算是涨了约 5.4\%，本章用全体城市消费者的指数，表 B-62 印的是 5.5\%。日本 1974 年的消费物价，不含自有住房估算房租的指数涨 24.5\%，综合指数涨 23.2\%，正文与图用前者。日本 1973 年的失业率，表里有不含冲绳与含冲绳两个数，1.2\% 与 1.3\%，图用后者。弗里德曼的正文说最后两个五年是 1966—1970 年与 1970—1975 年，他的表 1 写作 1971—1975 年。企划厅的网页版有个别排印的错字，如“基づものく”，引文照原样核对。
\end{deeper}
% 据: f68_scan, ph_title, mp_title, gk_title, b40_title, b62_title, lfs_table, cpi_xr_table, wp_title, gk_fn1, b62_1969, cpi_all_table, f76_one, wp_354

\section{今天还剩什么}\label{sec:naturalrate-legacy}

\textbf{自然率与中性利率。}弗里德曼 1976 年说，“自然失业率”这个说法，是他为了与维克塞尔的“自然利率”对应而造的\cite{friedman1976}。两个自然都指实际的力量决定、货币改变不了的那个水平，一个在劳动市场，一个在借贷市场（本书的解读）。今天的中央银行讲的中性利率，是维克塞尔的自然利率的后代（\chref{ch:wicksell}），泰勒规则按通胀与产出缺口调整利率（\chref{ch:consensus}）。
% 据: f76_wicksell

\textbf{预期。}本章的预期只看上一年的通胀。人们若会预期政策本身，结论还要更进一步，这是理性预期与卢卡斯批判（\chref{ch:ratex}）。降通胀的代价能不能靠可信的承诺减小，在\chref{ch:credibility}。名义刚性的新凯恩斯版本的曲线，在\chref{ch:newkeynesian}。

\textbf{失业的滞后与制度。}自然率若随实际的失业移动，长期的高失业就会自己延续下去，这种滞后效应与欧洲的失业是\chref{ch:labor}的题目。弗里德曼说，最低工资、工会与失业保险都会影响自然率\cite{friedman1968,friedman1976}。
% 据: f68_minimum, f76_insurance

\begin{takeaways}
\begin{itemize}
\item 弗里德曼 1967 年 12 月说，货币政策钉不住失业率。失业低于自然率，只能靠加速的通胀维持；通胀与失业之间只有暂时的取舍，没有永久的取舍。
\item 自然失业率是与实际工资结构的均衡相容的失业率。弗里德曼说它不是一成不变的，许多决定因素是人为的、政策造成的。
\item 附加预期的菲利普斯曲线是$\pi_t=\pi^{e}_t-\alpha(u_t-u^{n})$。在适应性预期下，失业低于自然率，通胀逐年上升；长期的曲线是$u^{n}$处的竖线。
\item 通胀降低多少，累计多出的失业就要多少，与快慢无关；奥肯估计牺牲率平均是 10，Gordon 与 King 认为只有约一半。代价更重地落在青年与黑人等本来就最容易失业的人身上（本书的概括）。
\item 1970—1985 年美国的通胀与失业几乎不再相关，通胀的变动却仍与失业反向。日本 1974 年物价涨了约四分之一，失业只有 1.4\%。
\end{itemize}
\end{takeaways}

\begin{selfcheck}
\begin{enumerate}
\item 按\exref{ex:naturalrate-accel}，把失业守在 4\%，前三年的通胀各是多少？
\item 按\exref{ex:naturalrate-disinflation}，通胀从 10\% 降到 4\%，失业守在 9\%，要几年？累计多出的失业是多少？
\item $\alpha$为多大时，按奥肯的系数 3.2 算出的牺牲率是 5？
\item 用本章的回归，失业率是 8\% 的那一年，通胀约变动多少？
\item 弗里德曼的表 1 里，美国从 1961—1965 年到 1966—1970 年，通胀与失业怎样变动？这符合简单的菲利普斯曲线吗？
\item 菲尔普斯为什么说他与弗里德曼发现的是两种不同的现象？
\end{enumerate}
\end{selfcheck}

\begin{answers}
\item 每年上升$0.5\times(6-4)=1$个百分点，依次是 3\%、4\% 与 5\%。
\item 每年降$0.5\times3=1.5$个百分点，要 4 年；累计多出 12 个百分点·年。
\item $3.2/\alpha=5$，$\alpha=0.64$。
\item $4.64-0.74\times8\approx-1.3$，约降 1.3 个百分点。
\item 通胀从 1.3\% 升到 4.2\%，失业从 5.5\% 降到 3.9\%，两者反向变动，符合。
\item 他说弗里德曼的模型讲的是劳动参与的自然率，他的讲的是失业的自然率；需求的意外增加，在前者是对竞争均衡的偏离，在后者减轻了非自愿的失业。
\end{answers}

回到 1967 年 12 月 29 日的华盛顿。那天晚上，弗里德曼讲完货币政策做不到的事，接着讲它能做的事\cite{friedman1968}。他的处方是货币当局公开宣布，让某一个货币总量稳定地增长\cite{friedman1968}。按他的估计，要让最终产品的物价大体稳定，现金加全部商业银行存款每年约要增长 3\% 到 5\%\cite{friedman1968}。本章的模型只说失业在长期里回到自然率，通胀由预期与失业决定（本书的概括）。货币增长与通胀之间还隔着一层货币需求，它稳不稳定，决定这个处方灵不灵（本书的解读）。控制货币数量就能控制通胀吗？
% 据: f68_topic, f68_prescription, f68_rule, f68_3to5b, f68_3to5
